\documentclass[11pt]{article}
\usepackage[letterpaper,margin=1in]{geometry}
\usepackage[T1]{fontenc}
\usepackage[utf8]{inputenc}
\usepackage{lmodern,microtype}
\usepackage{amsmath,amssymb,mathtools,bm,slashed}
\usepackage{xcolor,enumitem,booktabs,array,needspace}
\usepackage[most]{tcolorbox}
\usepackage{fancyhdr,titlesec,parskip,hyperref}
\definecolor{qftblue}{HTML}{245A8D}
\definecolor{qftgreen}{HTML}{2F6B4F}
\definecolor{qftgray}{HTML}{F4F6F8}
\hypersetup{colorlinks=true,linkcolor=qftblue,urlcolor=qftblue,pdfauthor={Nima Lashkari}}
\pagestyle{fancy}\fancyhf{}
\lhead{PHYS 662: Quantum Field Theory I}\cfoot{\thepage}
\setlength{\headheight}{14pt}
\titleformat{\section}{\Large\bfseries\color{qftblue}}{\thesection.}{0.6em}{}
\titleformat{\subsection}{\large\bfseries\color{qftgreen}}{\thesubsection}{0.6em}{}
\setlist[enumerate]{leftmargin=2em,itemsep=0.25em,topsep=0.35em}
\setlist[itemize]{leftmargin=2em,itemsep=0.2em,topsep=0.35em}
\newcommand{\dd}{\mathrm d}
\newcommand{\ii}{\mathrm i}
\newcommand{\eps}{\varepsilon}
\newcommand{\tr}{\operatorname{tr}}
\newcommand{\Tr}{\operatorname{Tr}}
\newcommand{\diag}{\operatorname{diag}}
\newcommand{\Lag}{\mathcal L}
\newcommand{\D}{\mathcal D}
\newcommand{\vac}{\lvert0\rangle}
\newcommand{\vbra}{\langle0\rvert}
\newcommand{\Tord}{\mathrm T}
\newcommand{\ket}[1]{\lvert#1\rangle}
\newcommand{\bra}[1]{\langle#1\rvert}
\newtcolorbox{keyidea}[1][]{enhanced,breakable,colback=qftgray,colframe=qftblue,
 boxrule=0.7pt,arc=1.5mm,title={#1},fonttitle=\bfseries,coltitle=white}
\newtcolorbox{remarkbox}[1][]{enhanced,breakable,colback=white,colframe=qftgreen,
 boxrule=0.6pt,arc=1.5mm,title={#1},fonttitle=\bfseries,coltitle=white}
\newcommand{\intp}{\int\!\frac{\dd^4p}{(2\pi)^4}}
\newcommand{\MSbar}{\overline{\mathrm{MS}}}
\newcommand{\module}[1]{\clearpage\section{#1}}
\newcommand{\lecturefront}[4]{%
\rhead{Lecture #1}\hypersetup{pdftitle={PHYS 662 Lecture #1: #2}}
\begin{document}\pagenumbering{gobble}
\begin{titlepage}\centering\vspace*{1.05in}
{\scshape\Large PHYS 662 --- Quantum Field Theory I\par}
\vspace{0.35in}{\Huge\bfseries\color{qftblue}Lecture #1\par}
\vspace{0.2in}{\LARGE\bfseries #2\par}
\vspace{0.2in}{\large\color{qftgreen}#3\par}
\vspace{0.6in}{\large Nima Lashkari\par}{\large Fall 2026\par}
\vfill\begin{minipage}{0.86\textwidth}\small #4\par\medskip
Course page: \url{https://nimalashkari.com/quantum-field-theory-2026-part-i/}
\end{minipage}\vspace{0.5in}\end{titlepage}
\pagenumbering{roman}\setcounter{tocdepth}{1}\tableofcontents
\vfill\begin{remarkbox}[Place in the course]
Lectures 19--20 established free scalar and fermion path integrals.
Lecture 21 turns the scalar source expansion into Feynman rules.
Lecture 22 evaluates the first scalar loops and introduces renormalization.
Problem Set 7 connects the diagram counting to these calculations.
\end{remarkbox}\clearpage\pagenumbering{arabic}}
\lecturefront{22}{Scalar Field Theory at One Loop}
{Mass, coupling, and the first renormalized amplitudes}
{We evaluate the tadpole and the three four-point bubbles of massive
real scalar theory. A cutoff gives a first view of the mass correction;
dimensional regularization then provides a common language for
counterterms, running parameters, and the physical two-particle cut.}

\section{Which quantities are being renormalized?}
Continue with \(\eta=\diag(1,-1,-1,-1)\), \(\hbar=c=1\), and the
unbroken \(\mathbb Z_2\)-symmetric theory. The line is
\(D_F(p)=\ii/(p^2-m^2+\ii0)\) and the quartic vertex is
\(-\ii\lambda\). Work perturbatively at \(m^2>0\). Our two targets
are the pole of the two-point function and the connected \(2\to2\)
scattering amplitude through one loop.

The parameters in a regulated action are called \emph{bare}. They
depend on the regulator if measured quantities are held fixed. Write
that same action in terms of a renormalized field and parameters:
\begin{align}
 \Lag&=\frac12(\partial\phi)^2-\frac12m^2\phi^2
 -\frac{\mu^{2\eps}\lambda}{4!}\phi^4+\Lag_{\rm ct},\notag\\
 \Lag_{\rm ct}&=\frac{\delta Z_\phi}{2}(\partial\phi)^2
 -\frac{\delta m^2}{2}\phi^2
 -\frac{\mu^{2\eps}\delta\lambda}{4!}\phi^4-\delta\Lambda_{\rm vac}.
 \label{eq:counterLag}
\end{align}
For dimensional regularization \(d=4-2\eps\); in strictly four
dimensions set \(\eps=0\). The scale \(\mu\) makes \(\lambda\)
dimensionless. If \(\phi_0=Z_\phi^{1/2}\phi\), then
\begin{equation}
 Z_\phi=1+\delta Z_\phi,\qquad
 Z_\phi m_0^2=m^2+\delta m^2,\qquad
 Z_\phi^2\lambda_0=\mu^{2\eps}(\lambda+\delta\lambda).
 \label{eq:bare}
\end{equation}
Counterterms are interaction vertices in the perturbative expansion.
The two-point vertex is \(\ii\delta Z_\phi p^2-\ii\delta m^2\);
the four-point vertex is \(-\ii\mu^{2\eps}\delta\lambda\).
The vacuum term cancels from our normalized correlators.

Fix the self-energy convention by writing a 1PI insertion as
\(-\ii\Sigma_R(p^2)\), including counterterms. Summing insertions gives
\begin{align}
 G_2(p)&=D_F+D_F(-\ii\Sigma_R)D_F+\cdots
 =\frac{\ii}{p^2-m^2-\Sigma_R(p^2)+\ii0},\label{eq:Dyson}\\
 \Sigma_R(p^2)&=\Sigma_{\rm loop}(p^2)+\delta m^2-\delta Z_\phi p^2.
 \label{eq:SigmaR}
\end{align}
The geometric expression organizes the perturbation series; a
one-loop insertion does not determine all higher-loop terms.
\begin{keyidea}[A regulator and a renormalization condition do different jobs]
A regulator makes intermediate expressions meaningful. A renormalization
condition specifies how \(m\) and \(\lambda\) are defined. Physical
predictions agree between schemes after those parameters are matched.
The renormalized mass need not equal the propagator's pole mass.
\end{keyidea}

\module{The tadpole and a Euclidean momentum cutoff}
Lecture 21 gave the one-loop two-point diagram with weight \(1/2\):
\begin{equation}
 -\ii\Sigma_{\rm tad}=\frac{-\ii\lambda}{2}
 \int\frac{\dd^4k}{(2\pi)^4}\frac{\ii}{k^2-m^2+\ii0},
 \qquad
 \Sigma_{\rm tad}=\frac{\lambda}{2}\int\frac{\dd^4k_E}{(2\pi)^4}
 \frac{1}{k_E^2+m^2}.                            \label{eq:tad}
\end{equation}
In the second expression rotate \(k^0=\ii k_E^0\) using the Feynman
pole prescription. The factor \(\ii\) in the measure and the signs
from the denominator and propagator combine to give the positive
Euclidean integrand. There is no external momentum in the loop.

A rotationally invariant Euclidean cutoff \(k_E^2\leq\Lambda^2\)
makes the calculation elementary. The unit three-sphere has area
\(2\pi^2\), hence
\begin{align}
 I_1^\Lambda(m^2)&=\frac{1}{8\pi^2}\int_0^\Lambda
 \frac{k^3\dd k}{k^2+m^2}
 =\frac{1}{16\pi^2}\left[\Lambda^2-m^2
 \log\!\left(1+\frac{\Lambda^2}{m^2}\right)\right],\label{eq:cutoff}\\
 \Sigma_{\rm tad}^\Lambda&=\frac{\lambda}{32\pi^2}
 \left[\Lambda^2-m^2\log\frac{\Lambda^2}{m^2}
 +O(m^4/\Lambda^2)\right].                        \label{eq:cutoffSigma}
\end{align}
The integral has mass dimension two, as a mass-squared correction
must. The quadratic divergence comes from momenta much greater than
\(m\); changing the ultraviolet regulator changes its coefficient.

Without counterterms, the pole would satisfy
\(M^2=m_0^2+\Sigma_{\rm tad}^\Lambda(m_0^2)+O(\lambda_0^2)\).
Holding the measured \(M\) fixed therefore requires a
cutoff-dependent \(m_0\). Equivalently, in an on-shell mass scheme
one sets \(m=M\) in the free propagator and chooses
\begin{equation}
 \delta m^2_{\rm OS}=-\Sigma_{\rm tad}^\Lambda(M^2),
 \qquad \delta Z_{\phi,\rm OS}=0\quad\hbox{at one loop}.
 \label{eq:OSmass}
\end{equation}
This is a definition of the mass parameter. It does not assert that
quantum corrections are absent. The relation between that parameter
and the bare action carries the correction.

\begin{remarkbox}[Normal ordering and momentum shifts]
Normal ordering \(\phi^4\) relative to a chosen free vacuum subtracts
its one-vertex self-contractions. This amounts to a particular mass
and vacuum-energy subtraction; it does not remove the four-point
bubble or the need for renormalization at higher orders.
A sharp cutoff also changes under a shift of loop momentum. For the
bubble we will use dimensional regularization, where such shifts are
consistent with the regulated integral.
\end{remarkbox}

\module{Dimensional regularization from a Gaussian integral}
Instead of imposing a boundary in momentum space, evaluate integrals
where they converge and analytically continue their dependence on
\(d\). For \(\Delta>0\) and \(\alpha>d/2\), the Schwinger identity
\begin{equation}
 \frac{1}{(k_E^2+\Delta)^\alpha}
 =\frac{1}{\Gamma(\alpha)}\int_0^\infty\dd t\,
 t^{\alpha-1}e^{-t(k_E^2+\Delta)}
\end{equation}
reduces the momentum integral to a Gaussian. Since
\(\int\dd^d k_E\,e^{-t k_E^2}/(2\pi)^d=(4\pi t)^{-d/2}\),
\begin{equation}
 I_\alpha(\Delta)=\mu^{2\eps}\int\frac{\dd^d k_E}{(2\pi)^d}
 \frac{1}{(k_E^2+\Delta)^\alpha}
 =\frac{\mu^{2\eps}}{(4\pi)^{d/2}}
 \frac{\Gamma(\alpha-d/2)}{\Gamma(\alpha)}
 \Delta^{d/2-\alpha},\qquad d=4-2\eps.             \label{eq:master}
\end{equation}
The right-hand side defines the analytic continuation. Introduce
\begin{equation}
 \frac{1}{\bar\eps}\equiv\frac1\eps-\gamma_E+\log(4\pi),
 \quad \Gamma(\eps)=\frac1\eps-\gamma_E+O(\eps),
 \quad \Gamma(-1+\eps)=-\frac1\eps+\gamma_E-1+O(\eps).
 \label{eq:epsbar}
\end{equation}
Also \(a^\eps=1+\eps\log a+O(\eps^2)\). Expanding \eqref{eq:master}
for \(\alpha=1,2\) gives the two integrals needed in this lecture:
\begin{align}
 I_1(m^2)&=-\frac{m^2}{16\pi^2}
 \left[\frac1{\bar\eps}+1-\log\frac{m^2}{\mu^2}\right]+O(\eps),
 \label{eq:I1}\\
 I_2(\Delta)&=\frac{1}{16\pi^2}
 \left[\frac1{\bar\eps}-\log\frac{\Delta}{\mu^2}\right]+O(\eps).
 \label{eq:I2}
\end{align}
The sign of \(I_1\) sometimes looks surprising. For \(d\) near four
this is an analytically continued expression, not an ordinary
convergent integral of a positive function. It need not have the
positive sign of the cutoff integral. There is no \(\Lambda^2\) term
in this scheme, but the mass-sensitive logarithm agrees after
parameters are matched.

\begin{keyidea}[Our subtraction convention]
Modified minimal subtraction, \(\MSbar\), removes the combination
\(1/\bar\eps\) and no additional finite term. All formulas use
\(d=4-2\eps\). A convention \(d=4-\epsilon\) has pole coefficients
twice as large when written as coefficients of \(1/\epsilon\).
\end{keyidea}
A scaleless integral is set to zero in dimensional regularization.
This fact alone does not prove the absence of ultraviolet divergences:
in a massless problem ultraviolet and infrared poles may cancel in a
scaleless expression. Keeping \(m>0\) avoids that ambiguity here.

\module{The renormalized mass and the pole residue}
Multiplying \eqref{eq:I1} by \(\lambda/2\), the tadpole becomes
\begin{equation}
 \Sigma_{\rm tad}=-\frac{\lambda m^2}{32\pi^2}
 \left[\frac1{\bar\eps}+1-\log\frac{m^2}{\mu^2}\right].
\end{equation}
It is independent of momentum, so it requires only a mass counterterm
at this order. In \(\MSbar\),
\begin{equation}
 \delta m^2=\frac{\lambda m^2}{32\pi^2\bar\eps}+O(\lambda^2),
 \qquad \delta Z_\phi=0+O(\lambda^2),\qquad
 \Sigma_R=\frac{\lambda m^2}{32\pi^2}
 \left[\log\frac{m^2}{\mu^2}-1\right]+O(\lambda^2).
 \label{eq:renmass}
\end{equation}
There is a useful direct sign check: the counterterm vertex is
\(-\ii\delta m^2\), so it adds \(+\delta m^2\) to our \(\Sigma\).
It cancels the negative pole in \(\Sigma_{\rm tad}\).

The pole and residue follow from the inverse propagator, not from
reading off the coefficient \(m\) in the Lagrangian:
\begin{align}
 M^2-m^2-\Sigma_R(M^2)&=0,                         \label{eq:polemass}\\
 Z_{\rm pole}&=\frac{1}{1-\Sigma_R'(M^2)}.         \label{eq:residue}
\end{align}
The prime differentiates with respect to \(p^2\). To first order,
\begin{equation}
 M^2=m^2(\mu)+\frac{\lambda(\mu)m^2(\mu)}{32\pi^2}
 \left[\log\frac{m^2(\mu)}{\mu^2}-1\right]+O(\lambda^2),
 \qquad Z_{\rm pole}=1+O(\lambda^2).              \label{eq:Mmass}
\end{equation}
Neither the scale-dependent \(\MSbar\) mass nor \(Z_\phi\) should be
confused with the measured pole mass or the residue of a chosen
renormalized field. Here both wavefunction corrections vanish at one
loop, but that coincidence does not persist.

If one instead defines \(m=M\), the on-shell counterterm in dimensional
regularization includes a finite part:
\begin{equation}
 \delta m^2_{\rm OS}=\frac{\lambda M^2}{32\pi^2}
 \left[\frac1{\bar\eps}+1-\log\frac{M^2}{\mu^2}\right]
 +O(\lambda^2).                                  \label{eq:OSdimreg}
\end{equation}
Then the renormalized one-loop self-energy is zero. Equations
\eqref{eq:renmass} and \eqref{eq:OSdimreg} describe different definitions
of the mass parameter, with the same pole after matching.

\begin{remarkbox}[What can be neglected in the next calculation?]
The one-loop four-point amplitude is already order \(\lambda^2\).
Replacing its internal \(m^2\) by \(M^2=m^2+O(\lambda)\) changes it
only at higher order, away from nonuniform threshold expansions.
The external kinematics must use \(p_a^2=M^2\).
Since the tree amplitude is order \(\lambda\), its residue factor
\(Z_{\rm pole}^2\) first matters at order \(\lambda^3\).
\end{remarkbox}

\module{Evaluating the four-point bubble}
Write \(\langle f|S-1|i\rangle=\ii(2\pi)^4\delta^{(4)}(P_f-P_i)
\mathcal M\), and define \(s,t,u\) as in Lecture 21. Each of the three
1PI bubbles has weight \(1/2\). Factor out the overall
\(\mu^{2\eps}\) of a four-point vertex before quoting its
four-dimensional amplitude. One channel then contributes
\begin{equation}
 \ii\mathcal M_q^{(1)}=\frac{\lambda^2}{2}\,J(q^2),\qquad
 J(q^2)=\mu^{2\eps}\int\frac{\dd^dk}{(2\pi)^d}
 \frac{1}{(k^2-m^2+\ii0)((k+q)^2-m^2+\ii0)}.
 \label{eq:bubbleintegral}
\end{equation}
The positive prefactor comes from the two vertices and the two
propagator numerators: \((-\ii\lambda)^2\ii^2=+\lambda^2\).
The integral itself supplies another factor of \(\ii\).

For positive Euclidean denominators,
\begin{equation}
 \frac1{AB}=\int_0^1\frac{\dd x}{[xA+(1-x)B]^2}.
\end{equation}
Take \(A=k_E^2+m^2\), \(B=(k_E+q_E)^2+m^2\). Completing the square
and shifting \(\ell_E=k_E+(1-x)q_E\) yields
\begin{equation}
 xA+(1-x)B=\ell_E^2+m^2+x(1-x)q_E^2.
\end{equation}
Use \eqref{eq:master} with \(\alpha=2\), then analytically continue
\(q_E^2\to-q^2-\ii0\). With
\(\Delta_x(q^2)=m^2-x(1-x)q^2-\ii0\),
\begin{align}
 J(q^2)&=\frac{\ii\mu^{2\eps}\Gamma(\eps)}{(4\pi)^{2-\eps}}
 \int_0^1\dd x\,[\Delta_x(q^2)]^{-\eps}\notag\\
 &=\frac{\ii}{16\pi^2}\left[\frac1{\bar\eps}-F(q^2;\mu)\right]
 +O(\eps),\qquad
 F(q^2;\mu)=\int_0^1\dd x\log\frac{\Delta_x(q^2)}{\mu^2}.
 \label{eq:bubbleanswer}
\end{align}
The logarithm is on the branch fixed by \(-\ii0\). This prescription
cannot be dropped when returning to physical timelike momenta.

Including tree, all three channels, and the local counterterm gives
\begin{equation}
 \mathcal M=-\lambda-\delta\lambda+
 \frac{\lambda^2}{32\pi^2}
 \left[\frac3{\bar\eps}-F(s;\mu)-F(t;\mu)-F(u;\mu)\right]
 +O(\lambda^3).                                  \label{eq:beforect}
\end{equation}
The pole is independent of \(s,t,u\). Its locality is why a single
\(\phi^4\) counterterm can cancel it for every scattering angle and
energy. The nonpolynomial momentum dependence in the finite logarithms
is a prediction, not a new independent coupling at each momentum.

\module{Coupling renormalization and a physical subtraction}
The \(\MSbar\) coupling counterterm and finite amplitude are
\begin{align}
 \delta\lambda&=\frac{3\lambda^2}{32\pi^2\bar\eps}+O(\lambda^3),
 \label{eq:deltalambda}\\
 \mathcal M(s,t,u)&=-\lambda(\mu)
 -\frac{\lambda(\mu)^2}{32\pi^2}
 \sum_{q^2=s,t,u}\int_0^1\dd x\,
 \log\frac{m^2-x(1-x)q^2-\ii0}{\mu^2}+O(\lambda^3).
 \label{eq:MSamplitude}
\end{align}
At this order the internal mass can be replaced by the pole mass
\(M\) as specified in Section 4, and physical kinematics obey
\(s+t+u=4M^2\). The equality of the three channel coefficients
makes the crossing symmetry of a real scalar manifest.

To see what a coupling measures, define a threshold coupling \(g\)
by \(\mathcal M(4M^2,0,0)=-g\), approached from the physical region.
At threshold the imaginary part vanishes. The elementary integral
\begin{equation}
 \int_0^1\dd x\log[1-4x(1-x)]=-2
\end{equation}
follows by writing \(1-4x(1-x)=(2x-1)^2\); the logarithmic singularity
at \(x=1/2\) is integrable. Hence
\begin{equation}
 g=\lambda(\mu)+\frac{\lambda(\mu)^2}{32\pi^2}
 \left[3\log\frac{M^2}{\mu^2}-2\right]+O(\lambda^3).
 \label{eq:matching}
\end{equation}
This finite conversion is required before comparing numerical values
of couplings in different schemes. Substituting it back gives
\begin{equation}
 \begin{aligned}
 \mathcal M(s,t,u)&=-g-\frac{g^2}{32\pi^2}
 \left[\sum_{q^2=s,t,u}f(q^2)+2\right]+O(g^3),\\
 f(q^2)&=\int_0^1\dd x\log\frac{M^2-x(1-x)q^2-\ii0}{M^2}.
 \end{aligned}
 \label{eq:thresholdamplitude}
\end{equation}
The bracket is zero at threshold. Neither the regulator nor \(\mu\)
appears. After two measurements fix \(M\) and \(g\), this expression
predicts a function of the scattering kinematics.

\begin{keyidea}[What is scheme dependent?]
Bare parameters, counterterms, and the numerical values of renormalized
parameters depend on the chosen definitions. A scattering amplitude
expressed in the same physical input data agrees between schemes,
up to terms beyond the perturbative order retained.
\end{keyidea}
The three bubble graphs are vertex corrections. Adding the four
external-leg tadpoles to \eqref{eq:MSamplitude} would double count the
pole and external-state treatment of Lecture 21. They belong to the
full correlator before LSZ reduction, not to this amputated scattering
formula.

\module{Scale independence and the one-loop beta function}
The arbitrary scale \(\mu\) cannot change a scattering experiment.
In \eqref{eq:MSamplitude}, each \(F\) has explicit derivative
\(\mu\partial_\mu F=-2\). At order \(\lambda^2\), implicit mass
running inside the loop contributes only at higher order. Therefore
\begin{equation}
 0=\mu\frac{\dd\mathcal M}{\dd\mu}
 =-\beta_\lambda+\frac{3\lambda^2}{16\pi^2}+O(\lambda^3),
 \qquad
 \boxed{\beta_\lambda\equiv\mu\frac{\dd\lambda}{\dd\mu}
 =\frac{3\lambda^2}{16\pi^2}+O(\lambda^3).}        \label{eq:beta}
\end{equation}
The derivative holds the bare theory fixed. This definition refers
to a renormalized coupling; it is not a derivative of a fixed bare
parameter with respect to \(\mu\).

The same result follows from
\(\lambda_0=\mu^{2\eps}(\lambda+a\lambda^2/\bar\eps+\cdots)\),
where \(a=3/(32\pi^2)\) and field renormalization enters this relation
only at higher order. Write
\(\beta_d=-2\eps\lambda+b\lambda^2+\cdots\). Differentiating gives
\begin{equation}
 0=2\eps\left(\lambda+\frac{a\lambda^2}{\bar\eps}\right)
 +\beta_d\left(1+\frac{2a\lambda}{\bar\eps}\right)+\cdots.
\end{equation}
The terms finite as \(\eps\to0\) imply \(b=2a\). Retaining the
classical \(-2\eps\lambda\) term is essential; dropping it too early
gives the wrong factor or sign.

Integrating the leading beta function between two scales gives
\begin{equation}
 \lambda(\mu)=\frac{\lambda(\mu_0)}
 {1-\dfrac{3\lambda(\mu_0)}{16\pi^2}\log(\mu/\mu_0)}.
 \label{eq:running}
\end{equation}
This resums leading logarithms. For positive coupling it grows toward
high energies. The formal denominator vanishes at
\(\mu_L=\mu_0\exp[16\pi^2/(3\lambda(\mu_0))]\), the one-loop Landau
scale. Perturbation theory fails before it can establish what happens
there; this formula is not a nonperturbative proof about the theory.

Scale independence also checks the mass calculation. Differentiating
\eqref{eq:Mmass} at fixed \(M\) gives
\begin{equation}
 \mu\frac{\dd m^2}{\dd\mu}=\frac{\lambda m^2}{16\pi^2}+O(\lambda^2).
\end{equation}
\begin{remarkbox}[Choosing a useful scale]
When kinematic invariants are of order \(Q^2\), choosing \(\mu\) of
order \(Q\) reduces large logarithms in the fixed-order amplitude.
In this mass-independent scheme, heavy particles do not automatically
decouple from the beta function; low-energy matching is a separate
step. The running parameter is a way of organizing predictions,
not an additional measurable dependence on an arbitrary scale.
\end{remarkbox}

\module{The two-particle threshold and an independent sign check}
The bubble logarithm develops an imaginary part when its two internal
particles can be on shell. In the physical \(s\) channel, define
\begin{equation}
 v=\sqrt{1-\frac{4M^2}{s}},\qquad x_\pm=\frac{1\pm v}{2}.
\end{equation}
For \(s>4M^2\), the argument \(M^2-sx(1-x)\) is negative precisely
for \(x_-<x<x_+\). On the Feynman branch
\(\log(-a-\ii0)=\log a-\ii\pi\) for \(a>0\), so
\begin{equation}
 \operatorname{Im}F(s;\mu)=-\pi v,\qquad
 \operatorname{Im}\mathcal M(s,t,u)
 =\frac{\lambda^2}{32\pi}\sqrt{1-\frac{4M^2}{s}}+O(\lambda^3).
 \label{eq:imag}
\end{equation}
In physical equal-mass \(2\to2\) scattering \(t,u\leq0\); their bubbles
have no such imaginary part. Local real counterterms cannot change
this discontinuity. Near threshold it is most useful to organize the
expansion with the physical mass in the propagators.

Unitarity gives an independent check of both the sign and the factor
\(1/2\). With \(S=1+\ii T\),
\begin{equation}
 -\ii(T-T^\dagger)=T^\dagger T.
\end{equation}
Insert a complete set of two-particle states into the forward matrix
element. At the first nontrivial order the optical theorem reads
\begin{equation}
 2\operatorname{Im}\mathcal M_{ii}^{(1)}
 =\frac1{2!}\int\dd\Phi_2\,
 |\mathcal M_{i\to2}^{\rm tree}|^2,
 \quad
 \dd\Phi_2=(2\pi)^4\delta^{(4)}(P-k_1-k_2)
 \prod_{a=1}^2\frac{\dd^3\bm k_a}{(2\pi)^3,2E_a}.
\end{equation}
In the center-of-mass frame the delta functions fix
\(|\bm k|=\sqrt{s/4-M^2}\), giving
\(\int\dd\Phi_2=v/(8\pi)\). Since
\(|\mathcal M^{\rm tree}|^2=\lambda^2\), the right-hand side is
\(\lambda^2v/(16\pi)\), exactly twice \eqref{eq:imag}.
The factorial avoids counting the same pair of identical intermediate
particles twice.

\begin{keyidea}[Results to carry forward]
The one-loop tadpole shifts the mass and has no momentum derivative.
The three bubbles require one local coupling counterterm. Their finite
parts depend on kinematics, their scale dependence gives
\(\beta_\lambda=3\lambda^2/(16\pi^2)\), and their branch cut describes
real two-particle intermediate states. Problem Set 7 checks these
statements from the source expansion through the optical theorem.
\end{keyidea}
{\footnotesize\textbf{Reading.} R.~G.~Leigh, \emph{Quantum Field Theory},
Sec.~6.5 and Sec.~6.7, for diagram organization and pole residues;
D.~Harlow, \emph{Quantum Field Theory I} (2024), Secs.~10.1 and 11,
for regulators, the tadpole, and the four-point bubble. All expressions
here use the conventions of this course.}
\end{document}
